% ============================================================
\chapter{Separation of Variables and Fourier Series}
\label{ch:fourier}
% ============================================================

In 1807, Joseph Fourier presented to the French Academy of Sciences a bold memoir: any periodic function, however irregular, can be written as a sum of sines and cosines. The claim was met with scepticism---Lagrange deemed it impossible---but it would prove to be one of the most fruitful ideas in all of mathematics. Fourier's method, born from the study of heat propagation, offers a remarkable strategy for solving linear PDEs: decompose the problem into eigenmodes, solve each mode separately, then recombine. This principle of \emph{separation of variables} transforms a PDE into a family of ordinary ODEs---a considerable gain.

\section{Principle of separation of variables}

The method of separation of variables (or Fourier's method) is one of the
most powerful techniques for solving linear PDEs on simple domains with
homogeneous boundary conditions.

\begin{intuition}
The idea is to look for the solution as a product of functions, each
depending on only one variable. The PDE then decomposes into several ODEs,
which are much easier to solve.
\end{intuition}

\subsection{General method}

Consider a PDE of the form $L[u] = 0$ posed on a domain
$\Omega = (0,L) \times (0,T)$ with homogeneous boundary conditions.

\textbf{Step 1:} Look for solutions of the form $u(x,t) = X(x)\, T(t)$.

\textbf{Step 2:} Substitute into the PDE and separate variables to obtain
two ODEs coupled by a separation constant $\lambda$.

\textbf{Step 3:} Solve the Sturm--Liouville problem for $X$ (with boundary
conditions), which determines the eigenvalues $\lambda_n$.

\textbf{Step 4:} Solve the ODE for $T$ with each $\lambda_n$.

\textbf{Step 5:} Form the general solution by superposition:
$u = \sum_n c_n X_n(x)\, T_n(t)$.

\textbf{Step 6:} Determine the coefficients $c_n$ from the initial
condition using Fourier series.

\begin{example}[Heat equation on a segment]
Consider:
\[
    \begin{cases}
        u_t = k\, u_{xx}, & 0 < x < L,\; t > 0,\\
        u(0,t) = u(L,t) = 0, & t > 0,\\
        u(x,0) = u_0(x), & 0 < x < L.
    \end{cases}
\]
Setting $u(x,t) = X(x)\, T(t)$, the PDE gives $X\, T' = k\, X''\, T$, i.e.:
\[
    \frac{T'}{kT} = \frac{X''}{X} = -\lambda \quad (\text{constant}).
\]
This yields two ODEs:
\[
    X'' + \lambda X = 0, \quad X(0) = X(L) = 0,
\]
\[
    T' + k\lambda T = 0.
\]
The problem for $X$ is a Sturm--Liouville problem whose nontrivial
solutions are:
\[
    \lambda_n = \left(\frac{n\pi}{L}\right)^2, \quad
    X_n(x) = \sin\!\left(\frac{n\pi x}{L}\right), \quad n \geq 1.
\]
The ODE for $T$ gives $T_n(t) = e^{-k\lambda_n t}$. The general solution is:
\[
    u(x,t) = \sum_{n=1}^{\infty} b_n \sin\!\left(\frac{n\pi x}{L}\right)
    e^{-k(n\pi/L)^2 t},
\]
where the $b_n$ are the Fourier sine coefficients of $u_0$.
\end{example}

\section{Fourier series}

\subsection{Definitions and coefficients}

\begin{definition}[Fourier series]
Let $f : [-\pi, \pi] \to \R$ be an integrable function. The \textbf{Fourier
series} of $f$ is:
\begin{equation}\label{eq:fourier_series}
    f(x) \sim \frac{a_0}{2} + \sum_{n=1}^{\infty}
    \left[a_n \cos(nx) + b_n \sin(nx)\right],
\end{equation}
where the \textbf{Fourier coefficients} are:
\begin{equation}\label{eq:fourier_coefficients}
    a_n = \frac{1}{\pi}\int_{-\pi}^{\pi} f(x)\cos(nx)\, dx, \quad
    b_n = \frac{1}{\pi}\int_{-\pi}^{\pi} f(x)\sin(nx)\, dx.
\end{equation}
\end{definition}

\begin{definition}[Complex Fourier series]
Equivalently, one can write:
\[
    f(x) \sim \sum_{n=-\infty}^{+\infty} c_n\, e^{inx}, \quad
    c_n = \frac{1}{2\pi}\int_{-\pi}^{\pi} f(x)\, e^{-inx}\, dx.
\]
\end{definition}

\begin{keyformulas}[title=\textbf{Fourier coefficients on $[0,L]$}]
For a function $f$ on $[0,L]$:
\begin{itemize}
    \item Sine coefficients: $b_n = \frac{2}{L}\int_0^L f(x)\sin\!\left(\frac{n\pi x}{L}\right) dx$.
    \item Cosine coefficients: $a_n = \frac{2}{L}\int_0^L f(x)\cos\!\left(\frac{n\pi x}{L}\right) dx$.
\end{itemize}
\end{keyformulas}

\subsection{Orthogonality}

\begin{proposition}[Orthogonality of trigonometric functions]
The system $\{1, \cos(nx), \sin(nx)\}_{n \geq 1}$ is orthogonal in
$L^2(-\pi, \pi)$:
\begin{align*}
    \inner{\cos(mx)}{\cos(nx)} &= \int_{-\pi}^{\pi} \cos(mx)\cos(nx)\, dx
    = \pi\, \delta_{mn}, \\
    \inner{\sin(mx)}{\sin(nx)} &= \int_{-\pi}^{\pi} \sin(mx)\sin(nx)\, dx
    = \pi\, \delta_{mn}, \\
    \inner{\cos(mx)}{\sin(nx)} &= 0 \quad \forall m, n.
\end{align*}
\end{proposition}

\subsection{Convergence of Fourier series}

\begin{theorem}[Pointwise convergence --- Dirichlet]
If $f$ is $2\pi$-periodic, piecewise continuous, and piecewise $C^1$,
then the Fourier series converges at every point:
\[
    \frac{a_0}{2} + \sum_{n=1}^{\infty} (a_n \cos(nx) + b_n \sin(nx))
    = \frac{f(x^+) + f(x^-)}{2}.
\]
In particular, at points of continuity the series converges to $f(x)$.
\end{theorem}

\begin{theorem}[Uniform convergence]
If $f$ is $2\pi$-periodic, continuous, and piecewise $C^1$, then the
Fourier series converges uniformly to $f$.
\end{theorem}

\begin{theorem}[Convergence in $L^2$ --- Parseval]\label{thm:parseval}
If $f \in L^2(-\pi, \pi)$, the Fourier series converges to $f$ in $L^2$
and \textbf{Parseval's identity} holds:
\begin{equation}\label{eq:parseval}
    \frac{1}{\pi}\int_{-\pi}^{\pi} \abs{f(x)}^2\, dx
    = \frac{a_0^2}{2} + \sum_{n=1}^{\infty} (a_n^2 + b_n^2).
\end{equation}
\end{theorem}

\begin{remark}
Parseval's identity expresses that the trigonometric system forms a
\textbf{Hilbert basis} of $L^2(-\pi,\pi)$. This is a special case of the
completeness theorem.
\end{remark}

\begin{example}
Let us compute the Fourier series of $f(x) = x$ on $[-\pi, \pi]$.
By oddness, $a_n = 0$ for all $n \geq 0$. For the $b_n$:
\[
    b_n = \frac{1}{\pi}\int_{-\pi}^{\pi} x \sin(nx)\, dx
    = \frac{2(-1)^{n+1}}{n}.
\]
Therefore:
\[
    x = 2\sum_{n=1}^{\infty} \frac{(-1)^{n+1}}{n}\sin(nx),
    \quad x \in (-\pi, \pi).
\]
Parseval's identity gives $\sum_{n=1}^{\infty} \frac{1}{n^2} = \frac{\pi^2}{6}$.
\end{example}

\section{Sturm--Liouville problems}

\begin{definition}[Sturm--Liouville problem]
A \textbf{regular Sturm--Liouville problem} is an eigenvalue problem of the
form:
\begin{equation}\label{eq:sturm_liouville}
    -(p(x)\, y')' + q(x)\, y = \lambda\, w(x)\, y, \quad x \in (a,b),
\end{equation}
with separated boundary conditions, where $p > 0$, $w > 0$, and $q$ are
continuous on $[a,b]$.
\end{definition}

\begin{theorem}[Properties of Sturm--Liouville problems]
Under the above hypotheses:
\begin{enumerate}
    \item There exists an infinite sequence of real eigenvalues
          $\lambda_1 < \lambda_2 < \cdots \to +\infty$.
    \item The corresponding eigenfunctions $\phi_n$ are orthogonal in
          $L^2_w(a,b)$:
          $\int_a^b \phi_m(x)\, \phi_n(x)\, w(x)\, dx = 0$ for $m \neq n$.
    \item The eigenfunction $\phi_n$ has exactly $n-1$ zeros in $(a,b)$.
    \item The system $\{\phi_n\}$ forms a Hilbert basis of $L^2_w(a,b)$.
\end{enumerate}
\end{theorem}

\begin{example}[Standard problems]
\leavevmode
\begin{enumerate}
    \item $X'' + \lambda X = 0$, $X(0) = X(L) = 0$ (Dirichlet):\\
          $\lambda_n = (n\pi/L)^2$, $X_n = \sin(n\pi x/L)$.
    \item $X'' + \lambda X = 0$, $X'(0) = X'(L) = 0$ (Neumann):\\
          $\lambda_n = (n\pi/L)^2$ for $n \geq 0$, $X_0 = 1$,
          $X_n = \cos(n\pi x/L)$.
    \item $X'' + \lambda X = 0$, $X(0) = 0$, $X'(L) = 0$ (mixed):\\
          $\lambda_n = ((2n-1)\pi/(2L))^2$, $X_n = \sin((2n-1)\pi x/(2L))$.
\end{enumerate}
\end{example}

\section{Applications to PDEs}

\subsection{Heat equation with Neumann conditions}

\begin{example}
Solve:
\[
    \begin{cases}
        u_t = k\, u_{xx}, & 0 < x < L,\; t > 0,\\
        u_x(0,t) = u_x(L,t) = 0, & t > 0,\\
        u(x,0) = u_0(x).
    \end{cases}
\]
Separation gives $X'' + \lambda X = 0$ with $X'(0) = X'(L) = 0$
(Neumann). The eigenvalues are $\lambda_n = (n\pi/L)^2$, $n \geq 0$, and:
\[
    u(x,t) = \frac{a_0}{2} + \sum_{n=1}^{\infty} a_n
    \cos\!\left(\frac{n\pi x}{L}\right) e^{-k(n\pi/L)^2 t},
\]
where $a_n = \frac{2}{L}\int_0^L u_0(x)\cos(n\pi x/L)\, dx$.

As $t \to \infty$, $u(x,t) \to a_0/2 = \frac{1}{L}\int_0^L u_0(x)\, dx$,
the initial mean temperature.
\end{example}

\subsection{Wave equation}

\begin{example}
Solve:
\[
    \begin{cases}
        u_{tt} = c^2 u_{xx}, & 0 < x < L,\; t > 0,\\
        u(0,t) = u(L,t) = 0,\\
        u(x,0) = u_0(x), \quad u_t(x,0) = v_0(x).
    \end{cases}
\]
Setting $u(x,t) = X(x)\, T(t)$ gives:
\[
    X'' + \lambda X = 0, \; X(0)=X(L)=0 \quad\text{and}\quad
    T'' + c^2\lambda T = 0.
\]
The solutions are $X_n = \sin(n\pi x/L)$, $\omega_n = cn\pi/L$, and:
\[
    u(x,t) = \sum_{n=1}^{\infty} \sin\!\left(\frac{n\pi x}{L}\right)
    \left[A_n \cos(\omega_n t) + B_n \sin(\omega_n t)\right],
\]
with:
\[
    A_n = \frac{2}{L}\int_0^L u_0(x)\sin\!\left(\frac{n\pi x}{L}\right) dx,
    \quad
    B_n = \frac{2}{L\omega_n}\int_0^L v_0(x)\sin\!\left(\frac{n\pi x}{L}\right) dx.
\]
\end{example}

\subsection{Laplace equation on a rectangle}

\begin{example}
Solve:
\[
    \begin{cases}
        \Delta u = 0, & 0 < x < a,\; 0 < y < b,\\
        u(0,y) = u(a,y) = 0,\\
        u(x,0) = 0, \quad u(x,b) = f(x).
    \end{cases}
\]
Setting $u = X(x)\, Y(y)$:
\[
    \frac{X''}{X} = -\frac{Y''}{Y} = -\lambda.
\]
With $X(0) = X(a) = 0$: $\lambda_n = (n\pi/a)^2$,
$X_n = \sin(n\pi x/a)$. The equation for $Y$ gives
$Y'' - \lambda_n Y = 0$ with $Y(0) = 0$:
\[
    Y_n(y) = \sinh\!\left(\frac{n\pi y}{a}\right).
\]
The solution is:
\[
    u(x,y) = \sum_{n=1}^{\infty} c_n \sin\!\left(\frac{n\pi x}{a}\right)
    \sinh\!\left(\frac{n\pi y}{a}\right),
\]
with $c_n = \frac{2}{a\sinh(n\pi b/a)}\int_0^a f(x)\sin(n\pi x/a)\, dx$.
\end{example}

\section{Separation in polar coordinates}

\begin{example}[Laplace in the disk]
Solve $\Delta u = 0$ in the disk $r < R$ with $u(R,\theta) = f(\theta)$.
In polar coordinates, $\Delta u = u_{rr} + \frac{1}{r}u_r + \frac{1}{r^2}u_{\theta\theta} = 0$.
Setting $u = \mathcal{R}(r)\, \Theta(\theta)$:
\[
    r^2\frac{\mathcal{R}''}{\mathcal{R}} + r\frac{\mathcal{R}'}{\mathcal{R}}
    = -\frac{\Theta''}{\Theta} = \lambda.
\]
Periodicity in $\theta$ forces $\lambda_n = n^2$, $n \geq 0$. The
solutions are:
\[
    u(r,\theta) = \frac{a_0}{2} + \sum_{n=1}^{\infty}
    \left(\frac{r}{R}\right)^n \bigl[a_n \cos(n\theta) + b_n \sin(n\theta)\bigr],
\]
where $a_n$, $b_n$ are the Fourier coefficients of $f$. This gives the
\textbf{Poisson integral formula}:
\[
    u(r,\theta) = \frac{1}{2\pi}\int_0^{2\pi}
    \frac{R^2 - r^2}{R^2 - 2Rr\cos(\theta-\phi) + r^2}\, f(\phi)\, d\phi.
\]
\end{example}

\section{Exercises}

\begin{exercise}
Compute the Fourier series of $f(x) = \abs{x}$ on $[-\pi,\pi]$.
Deduce that $\sum_{k=0}^{\infty} \frac{1}{(2k+1)^2} = \frac{\pi^2}{8}$.
\end{exercise}

\begin{exercise}
Solve the heat equation $u_t = u_{xx}$ on $(0,\pi)$ with
$u(0,t) = u(\pi,t) = 0$ and $u(x,0) = \sin(x) + 3\sin(2x)$.
\end{exercise}

\begin{exercise}
Solve the wave equation $u_{tt} = 4u_{xx}$ on $(0,1)$ with
$u(0,t) = u(1,t) = 0$, $u(x,0) = x(1-x)$, $u_t(x,0) = 0$.
\end{exercise}

\begin{exercise}
Solve the Laplace problem on the rectangle $(0,\pi) \times (0,1)$ with
$u(0,y) = u(\pi,y) = 0$, $u(x,0) = 0$, $u(x,1) = \sin(3x)$.
\end{exercise}

\begin{exercise}
Show that the Fourier coefficients of $f$ satisfy
$a_n, b_n \to 0$ as $n \to \infty$ (Riemann--Lebesgue lemma).
\end{exercise}

\begin{exercise}
Solve the Dirichlet problem in the disk $r < 1$ with
$u(1,\theta) = \cos^2\theta$. Write the solution in closed form.
\end{exercise}

\begin{exercise}[Singular Sturm--Liouville problem]
Find the eigenvalues and eigenfunctions of:
\[
    -\frac{d}{dx}\left(x\frac{dy}{dx}\right) = \frac{\lambda}{x}\, y,
    \quad 1 < x < e, \quad y(1) = y(e) = 0.
\]
\emph{Hint:} set $x = e^s$.
\end{exercise}
